% paper/sections/06c_fccd_advection.tex
%
% §6.2  FCCD 移流：Option B/C と面ジェット共通プリミティブ (CHK-223: §6 per-term spatial 配下)
%        既存の節点 CCD 移流を保存型に拡張する FCCD 移流演算子群．
%        §4.5 FCCD の面中心行列 M^FCCD を速度移流に転用する経路を示す．
%
% ═══════════════════════════════════════════════════════════════════════════════
\subsection{FCCD 面移流演算子：面ロカス閉包と参照再構成}
\label{sec:fccd_advection}
% ═══════════════════════════════════════════════════════════════════════════════

\noindent\textbf{Part 1 連続定式化との関係：}
本節の FCCD 面移流演算子は §~\ref{sec:onefluid}（§2 One-Fluid NS，
式~\eqref{eq:NS_onefluid}）の運動量移流項を，
Balanced--Force の同一面ロカス原理（§~\ref{sec:balanced_force}）と整合する形で
記述するための演算子契約である．本稿の標準 bulk 運動量移流は
UCCD6（§~\ref{sec:uccd6_def}）であり，FCCD 面フラックス形は
縮約 CSF 経路で圧力・毛管・対流フラックスを同一 face-to-node divergence で閉じる場合，
または標準 pressure-jump 経路で毛管 face cochain と動的対流フラックスのロカスを比較する場合の
面ロカス構成として用いる．純 FCCD DNS 構成での役割は §~\ref{sec:pure_fccd_dns} で境界づける．

本節では \S\ref{sec:fccd_def} で導入した面中心 FCCD 演算子
$\mathbf{M}^{\FCCD}$（式~\eqref{eq:fccd_composite_matrix}）を
速度移流項へ拡張し（保存型射影法での速度移流の標準解説は~\cite{BellColella1989}），
Balanced--Force 原理（\S\ref{sec:balanced_force}）を一括 CSF 縮約で読む場合に現れる
「$\bnabla p,\,\sigma\kappa_{lg}\bnabla\psi,\,\bnabla\cdot(\bu\otimes\bu)$
\textbf{3 項すべてが同一の面中心演算子}で離散化される」条件を完結させる．
標準 pressure-jump 経路では，この条件を毛管 face cochain と正準面補正の同一ロカス条件として読み替える．
節点版の運動量移流を圧力・毛管面ロカスと混在させると，
\textbf{移流だけが節点中心に残る}ため，静的 BF 条件は保てても
動的 RHS に節点--面の往復誤差が残る．
本節では 2 つの役割を区別する：
\begin{itemize}
  \item \textbf{節点再構成形（Option C）}：面中心 FCCD 勾配を R$_4$
        Hermite 再構成で節点に戻す参照構成．既存の節点 RHS と比較するときの
        精度分解に用いる．
  \item \textbf{保存形面フラックス（Option B）}：
        $\bnabla\cdot(\bu\otimes\bu)$ を FCCD 面値 $\mathbf{P}_f$ で
        組み立て，面ダイバージェンスで節点に戻す．動的 BF ロカス整合を
        最も直接に満たす構成である．
\end{itemize}

% -------------------------------------------------------------------------------
\subsubsection{\texorpdfstring{共通プリミティブ：面値演算子 $\mathbf{P}_f$ と面ジェット $\mathcal{J}_f$}{共通プリミティブ：面値演算子 P\_f と面ジェット J\_f}}
\label{sec:fccd_adv_primitives}
% -------------------------------------------------------------------------------

節点ベクトル $\mathbf{u}\in\mathbb{R}^{N+1}$ から面ベクトル
$\mathbf{u}_f\in\mathbb{R}^N$ への 4 次精度内挿は，
節点平均の $\mathcal{O}(H^2)$ 誤差を $\mathbf{S}_{\CCD}$ が与える第二導関数
$\mathbf{q}\in\mathbb{R}^{N+1}$ で相殺することで得られる．
面中央 $x_c = x_{f_{i-1/2}}$ 周りの対称 Taylor 展開から：
\begin{equation}
  u^{\FCCD}_{f_{i-1/2}}
  = \tfrac{1}{2}(u_{i-1}+u_i)
  - \tfrac{H_i^2}{16}(q_{i-1}+q_i)
  + \mathcal{O}(H^4).
  \label{eq:fccd_adv_face_value}
\end{equation}
疎な双対角行列 $\mathbf{P}_1, \mathbf{P}_2\in\mathbb{R}^{N\times(N+1)}$ を
$\mathbf{P}_1 = \tfrac{1}{2}(1,1)\text{ 双対角}$，
$\mathbf{P}_2 = (H_i^2/16)\cdot(1,1)\text{ 双対角}$ として，
\begin{equation}
  \boxed{\;
  \mathbf{P}_f := \mathbf{P}_1 - \mathbf{P}_2\,\mathbf{S}_{\CCD},
  \qquad \mathbf{u}_f = \mathbf{P}_f\,\mathbf{u},
  \;}
  \label{eq:fccd_adv_Pf}
\end{equation}
と行列化する．
切断誤差 $u_f - u(x_c) = -\tfrac{5}{384}H^4 u^{(4)}(x_c) + \mathcal{O}(H^6)$．

\noindent\textbf{面ジェット $\mathcal{J}_f$}（\S\ref{sec:face_jet_def}）：
統一プリミティブ $\FaceJet{u} = (u_f, u'_f, u''_f)$ は
\begin{equation}
  \mathcal{J}_f(\mathbf{u}) = \bigl(\mathbf{P}_f\mathbf{u},\,
  \mathbf{M}^{\FCCD}\mathbf{u},\,
  \mathbf{S}_{\CCD}^f\mathbf{u}\bigr),
  \label{eq:fccd_adv_jet}
\end{equation}
の 3 出力を単一の $\mathbf{S}_{\CCD}$ 解から共有する．
面値・面勾配・面第二導関数はいずれも同じ $\mathbf{q}$ に依存するため
\textbf{CCD ブロック三重対角解は 1 軸あたり 1 回のみ}で済む．

% -------------------------------------------------------------------------------
\subsubsection{\texorpdfstring{Option C：R$_4$ Hermite 面→節点再構成}{Option C：R\_4 Hermite 面→節点再構成}}
\label{sec:fccd_adv_option_c}
% -------------------------------------------------------------------------------

面値勾配 $\mathbf{d}^{\FCCD} = \mathbf{M}^{\FCCD}\mathbf{u}$ から
4 次精度の節点勾配を得るには，隣接面の平均で生じる
$\mathcal{O}(H^2)$ 誤差を $\mathbf{q}$ で相殺する：
\begin{equation}
  (\partial_x u)^{\text{node-FCCD}}_i
  = \tfrac{1}{2}(d_{f_{i-1/2}} + d_{f_{i+1/2}})
  - \tfrac{H}{16}(q_{i+1} - q_{i-1})
  + \mathcal{O}(H^4).
  \label{eq:fccd_adv_R4}
\end{equation}
この\textbf{R$_4$ Hermite 再構成}で節点勾配が FCCD の面精度と
同じ $\mathcal{O}(H^4)$ を達成する．
移流項は従来通り node-pointwise 積で組み立てる：
\begin{equation}
  [\mathcal{A}^{(j)}]_i = -\sum_k u_k(x_i)\cdot(\partial_{x_k} u_j)^{\text{node-FCCD}}_i.
  \label{eq:fccd_adv_option_c}
\end{equation}

\noindent\textbf{Option C の役割}：
出力が節点値なので，節点中心 RHS を保ったまま
FCCD 面勾配の影響だけを切り出して比較できる．
したがって Option C は節点中心経路との精度差を説明する参照構成であり，
標準の二相 face-closure を置き換えるものではない．

\noindent\textbf{Option C の限界}：
R$_4$ の面→節点変換で BF 整合性が\textbf{部分的に}回復するに留まる．
縮約 CSF 経路の圧力・毛管項が面中心 FCCD で動く一方，移流だけは R$_4$ 経由の
節点に戻っているため，節点-面-節点の往復で
$\mathcal{O}(H^4)$ 誤差が 2 重に蓄積する．
面ロカス閉包を主張する箇所では Option B を用いる．

% -------------------------------------------------------------------------------
\subsubsection{Option B：保存形 FCCD 面フラックス}
\label{sec:fccd_adv_option_b}
% -------------------------------------------------------------------------------

運動量の自然形は $\bnabla\cdot(\bu\otimes\bu)$（保存形）である．
非保存形 $(\bu\cdot\bnabla)\bu$ は $\bnabla\cdot\bu\neq 0$
の離散的残差によりスプリアス運動量源を生む．
FCCD 面フラックスで保存形を評価する：
\begin{equation}
  F^{(k,j),\text{cons}}_{f_{i-1/2}}
  = (u^{(k)} u^{(j)})_f
  = \bigl(\mathbf{P}_f(u^{(k)} u^{(j)})\bigr)_{i-1/2}
  \quad\text{（節点積 $u^{(k)}_i u^{(j)}_i$ を先に作り，単一の $\mathbf{P}_f$ で面に補間）},
  \label{eq:fccd_adv_flux_cons}
\end{equation}
面で 2 回分離補間する $(\mathbf{P}_f u^{(k)})(\mathbf{P}_f u^{(j)})$ は
非線形積の補間誤差が $\Ord{H^2}$ フロアに落ちるため採用しない．
節点ダイバージェンスで運動量更新を閉じる：
\begin{equation}
  [\bnabla\cdot F^{(j)}]_i
  = \sum_k\frac{F^{(k,j)}_{f_{i+1/2}} - F^{(k,j)}_{f_{i-1/2}}}{H_i}.
  \label{eq:fccd_adv_divergence}
\end{equation}

\noindent\textbf{スキュー対称形（推奨）}：
運動量の離散エネルギー保存のため，保存形と非保存形の算術平均
\begin{equation}
  F^{(k,j),\text{sk}}_{f_{i-1/2}}
  = \tfrac{1}{2}\bigl(F^{(k,j),\text{cons}}_{f_{i-1/2}}
  + F^{(k,j),\text{nc}}_{f_{i-1/2}}\bigr)
  \label{eq:fccd_adv_flux_skew}
\end{equation}
を用いる（Kok 2000, Ham et al. 2002）．
非保存形 $F^{(k,j),\text{nc}} = (\mathbf{P}_f u^{(k)})\cdot(\mathbf{M}^{\FCCD}u^{(j)})$
は面値と面勾配の点積，両者とも $\mathcal{J}_f$ から得られる．

\smallskip\noindent\textbf{離散エネルギー保存の面ジェット版（要約）}：
発散フリー速度場 $D_h\bu=0$ で，
$\langle\psi,\mathcal{L}^\text{sk}\psi\rangle = \tfrac{1}{2}\langle\psi,\mathcal{L}^\text{cons}\psi\rangle
+ \tfrac{1}{2}\langle\psi,\mathcal{L}^\text{nc}\psi\rangle$ より
保存形と非保存形の対称化により $\langle\psi,\mathcal{L}^\text{sk}\psi\rangle=0$
が \textbf{離散積} $\langle\cdot,\cdot\rangle_{f}$（面ジェット重み付き内積）で成立．
ここで CCD 面ロカス $\mathbf{P}_f$ と面勾配 $\mathbf{M}^{\FCCD}$ が
\textbf{随伴ペア}を構成することが鍵で，
$\langle\mathbf{P}_f\mathbf{u},\mathbf{M}^{\FCCD}\mathbf{v}\rangle_f
=\langle\mathbf{u},(\mathbf{M}^{\FCCD})^\top\mathbf{P}_f^\top\mathbf{v}\rangle$
の関係（一様格子では境界寄与を除き厳密；非一様格子では $\Ord{H^4}$
誤差で近似）が成立する．これは P-2 の補正幾何経路統一
（§\ref{sec:bf_p2_adjoint}）の帰結である．

% -------------------------------------------------------------------------------
\subsubsection{BF 整合性定理（面上完結）}
\label{sec:fccd_adv_bf_theorem}
% -------------------------------------------------------------------------------

\begin{theorem}[縮約 CSF 経路での FCCD 面ロカス静的 BF 整合性]\label{thm:fccd_adv_bf}
一括 CSF 経路で $\bnabla p$ と $\sigma\kappa_{lg}\bnabla\psi$ を
\textbf{共通の面中心演算子}
$\mathbf{M}^{\FCCD}$（\S\ref{sec:fccd_def} 式~\eqref{eq:fccd_composite_matrix}）で
離散化すれば，静止状態 $\bu\equiv 0$ での離散運動量バランスの
BF 残差は
\begin{equation*}
  \|\mathbf{R}_\text{BF}\|_\infty = \mathcal{O}(H^4)\text{（一様）},
  \qquad \mathcal{O}(H^3)\text{（非一様）}
\end{equation*}
を満たす．
\end{theorem}

\begin{proof}[\upshape 証明（要旨）]
\textbf{仮定の明示}：(A1) FCCD 共通演算子 $\mathbf{M}^{\FCCD}$ は
\S\ref{sec:fccd_def}（一様）で切断誤差 $\Ord{H^4}$，
\S\ref{sec:fccd_nonuniform}（非一様）で $\Ord{H^3}$．
(A2) PPE 可解条件 $\sum_i (D_h\bu^*)_i = 0$（離散発散の総和零；
Neumann BC 下の compatibility）が満たされる．
(A3) 本定理は\textbf{一括 CSF PPE}（変密度一括 PPE + smoothed Heaviside
CSF；§\ref{sec:projection_derivation}）の設定．分相 affine-jump PPE 版
（§\ref{sec:split_ppe_pressure_jump_formulation}）は別建ての
Corollary~\ref{cor:fccd_split_ppe_bf}．\\
\textbf{論証}：$\bu\equiv 0$ では対流項は恒等的に零であるため，
静的 BF 残差は圧力勾配と縮約 CSF 毛管項の演算子一致だけで決まる．
残る $\bnabla p$ と縮約 CSF 毛管項は (A1) より同一の
$\mathbf{M}^{\FCCD}$ で評価され，両項の切断誤差は同一表現
$\Ord{H^4}/\Ord{H^3}$．差分 $G_h p - f_\sigma$ では主項が打ち消し合い
（同一演算子・同一格子点・同一切断），$\|\mathbf{R}_\text{BF}\|_\infty$ は
切断誤差オーダーで律速される．(A2) は PPE 可解性を保証し，$p^{n+1}$ が
well-defined であることから $G_h p$ の評価が有効．$\square$
\end{proof}

\begin{corollary}[分相 affine-jump PPE での FCCD-BF 整合]
\label{cor:fccd_split_ppe_bf}
分相 affine-jump PPE（§\ref{sec:split_ppe_framework}）を採用した場合，
CSF 体積力は予測段階から外れ，毛管効果は
$J_p^\chi=-\sigma\kappa_{lg}$ の face cochain として PPE/corrector に入る．
このとき BF 整合性は，HFE が供給する片側 Hermite データを
§\ref{sec:split_ppe_pressure_jump_formulation} のジャンプ補正付き面法則へ渡し，
さらに §\ref{sec:split_ppe_capillary_range_projection} の
$D_f A_f G_f$ 値域射影後の cochain を正準面補正へ渡すことで成立する．
節点延長圧力の raw 勾配を速度補正へ直接渡す経路はこの corollary の仮定ではない．
\end{corollary}

\noindent\textbf{動的移流ロカス命題}：
非零速度場では，対流フラックスも圧力・毛管 cochain と同じ面ロカスに置かれて初めて
運動量 RHS 全体が同一の面中心空間で閉じる．
Option C（R$_4$ 経由の節点中心）は節点--面--節点の往復を残すため，
静的 BF 定理の直接の仮定ではないが，動的計算ではロカス不一致の残差源になり得る．
Option B は $\bnabla\cdot(\bu\otimes\bu)$ を面フラックスとして構成し，
縮約 CSF 経路では圧力・毛管・対流の 3 項を，標準 pressure-jump 経路では
毛管 face cochain と対流 face flux の比較基準を，同一の face-to-node divergence で読むための設計である．

% -------------------------------------------------------------------------------
\subsubsection{運動量方程式の閉形}
\label{sec:fccd_adv_closure}
% -------------------------------------------------------------------------------

FCCD Option B で閉じた運動量更新：
\begin{equation}
  \frac{u^{n+1}_j - u^n_j}{\Delta t}
  = -\sum_k\frac{(u_k u_j)^{*}_{f_{i+1/2}} - (u_k u_j)^{*}_{f_{i-1/2}}}{H_i}
  + \bigl[\mathcal{V}(\bu) - \bnabla p + \sigma\kappa_{lg}\bnabla\psi + \rho\bm{g}\bigr]_i,
  \label{eq:fccd_adv_closure}
\end{equation}
ここで $\mathcal{V}(\bu)=\bnabla\cdot[\mu(\bnabla\bu+\bnabla\bu^T)]$ は
§\ref{sec:onefluid}（§2 NS 既定形）と §\ref{sec:viscous_3layer}（§6.3 3 層粘性）
で定義された粘性演算子．この式は縮約 CSF 経路の面ロカスを示す代表形であり，
右辺は粘性・圧力・縮約毛管項・重力を\textbf{面中心}で評価し
面ダイバージェンスで節点に戻す．
PPE RHS は自然に $\bnabla_h\cdot\bu^*$ が FCCD 面ダイバージェンスで
表現され，ループが閉じる．

% -------------------------------------------------------------------------------
\subsubsection{壁境界条件 Option IV：Dirichlet 無滑り}
\label{sec:fccd_adv_wall_iv}
% -------------------------------------------------------------------------------

\S\ref{sec:ccd_bc} で導入した Option III（Neumann $\psi, p$）は
移流でも変更不要．
Dirichlet 無滑り $u(x_\text{wall}) = 0$ に対する
\textbf{Option IV（Dirichlet 壁境界）}は，
$u$ の反対称ゴーストミラー $u_{-1} = -u_1$ と
$u''$ の反対称性 $q_{-1} = -q_1$ を組み合わせる：
\begin{align}
  d^{\FCCD}_{f_{-1/2}} &= \frac{u_1}{H} - \frac{H}{24}(q_0 + q_1)
  \quad\text{（壁せん断率, 一般に非零）}, \label{eq:fccd_adv_wall_iv_d}\\
  u^{\FCCD}_{f_{-1/2}} &= -\tfrac{1}{2}u_1 + \tfrac{H^2}{16}(q_1 - q_0)
  \quad\text{（壁面速度, 連続無滑り極限で零）}. \label{eq:fccd_adv_wall_iv_u}
\end{align}
Option B の壁面フラックス
$F^{(k,j)}_{f_{-1/2}} = u^{(k)}_f\cdot u^{(j)}_f$ は
両因子とも壁で $\to 0$ のため先頭オーダーで恒等零，
物理的無滑り運動量フラックス零と一致．

% -------------------------------------------------------------------------------
\subsubsection{演算量・役割分担}
\label{sec:fccd_adv_impl}
% -------------------------------------------------------------------------------

\noindent\textbf{$\mathbf{q}$-共有最適化}：
面値 $\mathbf{P}_f\mathbf{u}$，面勾配 $\mathbf{M}^{\FCCD}\mathbf{u}$，
面第二導関数 $\mathbf{S}_{\CCD}^f\mathbf{u}$ はすべて共通の
$\mathbf{q} = \mathbf{S}_{\CCD}\mathbf{u}$ に依存するため，
CCD ブロック三重対角解は\textbf{1 軸あたり 1 回のみ}．
追加コストは疎行列 $\mathbf{P}_1, \mathbf{P}_2, \mathbf{D}_1, \mathbf{D}_2$ の
$\mathcal{O}(N)$ 乗算のみで，
演算量は純 CCD 移流の 1.3--1.5$\times$．

\noindent\textbf{本稿内での役割分担}：
\begin{itemize}
  \item 標準 bulk 運動量移流は UCCD6（§\ref{sec:uccd6_def}）で担う．
  \item 一括 CSF 縮約経路の BF 説明では，共通面ロカスとして
        圧力・毛管・対流フラックスの整合条件を本節の Option B で表現できる．
  \item 分相 affine-jump 経路では，毛管効果は CSF 体積力ではなく
        §\ref{sec:split_ppe_pressure_jump_formulation} の face cochain として入り，
        capillary range projection 後に正準面補正へ渡る．
  \item Option C は節点中心 RHS と面中心 RHS の差を切り分ける参照構成であり，
        標準 face-closure の代替ではない．
\end{itemize}

% ═══════════════════════════════════════════════════════════════════════════════
